\documentclass[8pt]{beamer}

% Beamer style
\useoutertheme{infolines}
\usecolortheme[rgb={0.65,0.15,0.25}]{structure}
\beamertemplatenavigationsymbolsempty

% Packages
\usepackage[french]{babel}
\usepackage[latin1]{inputenc}
\usepackage{color}
\usepackage{xspace}
\usepackage{dsfont, stmaryrd}
\usepackage{amsmath, amsfonts, amssymb, stmaryrd}
\usepackage{epsfig}
\usepackage{tikz}
\usepackage{url}
\usepackage{/home/robin/LATEX/Biblio/astats}
\usepackage{graphicx}
% \usepackage{enumitem}

\input{/home/robin/RECHERCHE/EXPOSES/LATEX/SlideCommands}
\newcommand{\GMSBM}{/home/robin/RECHERCHE/RESEAUX/EXPOSES/1903-SemStat/}
\newcommand{\figeconet}{/home/robin/RECHERCHE/ECOLOGIE/EXPOSES/1904-EcoNet-Lyon/Figs}
\newcommand{\fignet}{/home/robin/RECHERCHE/RESEAUX/EXPOSES/FIGURES}
\newcommand{\figeco}{/home/robin/RECHERCHE/ECOLOGIE/EXPOSES/FIGURES}
\newcommand{\figbayes}{/home/robin/RECHERCHE/BAYES/EXPOSES/FIGURES}
\newcommand{\figCMR}{/home/robin/Bureau/RECHERCHE/ECOLOGIE/CountPCA/sparsepca/Article/Network_JCGS/trunk/figs}
\newcommand{\figtree}{/home/robin/RECHERCHE/BAYES/VBEM-IS/VBEM-IS.git/Data/Tree/Fig}
\newcommand{\figDoR}{/home/robin/RECHERCHE/BAYES/VBEM-IS/VBEM-IS.git/Paper/JRSSC-V3/Figs}

%====================================================================
%====================================================================
\begin{document}
%====================================================================
%====================================================================

%====================================================================
%====================================================================
\section*{Backup}
%====================================================================
\frame{\frametitle{L3 : ``Ce qu'un biologiste doit savoir en math�matiques''}

  \paragraph{Cours existant.} Conjoint avec l'ENS Ulm

  \bigskip \bigskip 
  \paragraph{Programme.} \\ ~
  \begin{itemize}
   \setlength{\itemsep}{1.2\baselineskip}
   \item Alg�bre lin�aire
   \item Diff�rentiabilit�
   \item \'Equations diff�rentielles
   \item Probabilit�s
  \end{itemize}
  
}

%====================================================================
\frame{\frametitle{L1 : ``Sciences des donn�es''}

  \paragraph{Cours � construire.} 
  Conjoint avec l'UFR d'informatique

  \bigskip \pause
  \paragraph{Parti pris.} 
  \begin{itemize}
   \item Importance des m�thodes computationnelles dans l'analyse des donn�es
   \item Distinction entre analyse descriptive et inf�rence
   \item D�monstration des propri�t�s de m�thodes pr�sent�es
   \item Pas d'introduction de nouvelle notion math�matique 
   \item Impl�mentation des m�thodes en TD
  \end{itemize}

  \bigskip \pause
  \paragraph{Quelques exemples} s'appuyant sur le programme de terminale 
  \begin{itemize}
   \item R�gression lin�aire $\neq$ composante principale (dimension 2)
   \item Classification non-supervis�e et algorithme des $k$-means
   \item D�tection de ruptures et programmation dynamique
   \item M�thode de Newton pour l'obtention d'un estimateur par optimisation d'un crit�re
   \item Estimation d'une int�grale par Monte-Carlo
   \item Validation crois�e et sur-ajustement
  \end{itemize}
  
}

%====================================================================
%====================================================================
\end{document}
%====================================================================
%====================================================================

  \begin{tabular}{cc}
    \hspace{-.04\textwidth}
    \begin{tabular}{p{.5\textwidth}}
    \end{tabular}
    &
    \begin{tabular}{p{.45\textwidth}}
    \end{tabular}
  \end{tabular}

